% ============================================================================
%  DOCUMENTO PRINCIPAL
%  ----------------------------------------------------------------------------
%  Aquí SOLO se integran los módulos. No se escribe contenido ni formato aquí.
%
%  Regla de oro del proyecto:
%    - formato/preambulo.sty  -> toda la configuración y paquetes (APA + Arial)
%    - formato/portada.tex    -> portada (datos institucionales)
%    - contenido/*.tex        -> el contenido real del informe
%    - anexos/*.tex           -> material complementario
%
%  Para compilar (desde latex/):
%    make
% ============================================================================

\documentclass[12pt,a4paper]{article}

% --- Carga de la configuración de formato (APA + Arial) ---------------------
\usepackage{formato/preambulo}

% ============================================================================
%  DATOS DEL DOCUMENTO  (editar aquí)
% ============================================================================
\newcommand{\tituloDocumento}{Diseño e Implementación de Infraestructura de Red GNU/Linux para el Banco Sudoers}
\newcommand{\equipo}{Grupo 1 Sudoers}
\newcommand{\integrantes}{%
  Joaquin Alessandro Felipez Rojas \par
  Nicolas Sebastian Reguerin Meneses \par
  David Willy Cruz Huanca}
\newcommand{\universidad}{Universidad Privada Domingo Savio}
\newcommand{\facultad}{Facultad de Ingenier\'ia}
\newcommand{\carrera}{Ingenier\'ia en Sistemas}
\newcommand{\asignatura}{Taller de Sistemas Operativos II}
\newcommand{\docente}{Jose Abraham Mita Huanca}
\newcommand{\ciudad}{Cochabamba--Bolivia}
\newcommand{\fechaEntrega}{28 de septiembre de 2026}

\begin{document}

% --- Numeración: romano antes del Capítulo I, arábigo desde el Capítulo I ---
\pagenumbering{roman}

% --- Portada (no muestra número y no cuenta) ---
% ============================================================================
%  PORTADA  (formato APA 7.ª edición)
%  ----------------------------------------------------------------------------
%  Los datos provienen de las macros definidas en main.tex. Editar ahí.
% ============================================================================

\begin{titlepage}
  \thispagestyle{empty}
  \setstretch{1}
  \setlength{\parskip}{0pt}

  \begin{center}
    \vspace*{0.6cm}

    {\bfseries\fontsize{20}{24}\selectfont \universidad \par}
    {\bfseries\fontsize{18}{22}\selectfont \facultad \par}
    {\bfseries\fontsize{16}{19}\selectfont \carrera \par}

    \vspace{0.8cm}

    % Logo de la institución
    \IfFileExists{figuras/logo.png}{%
      \includegraphics[width=0.30\textwidth]{figuras/logo.png}\par
    }{}

    \vspace{0.9cm}

    {\bfseries\Large \tituloDocumento \par}

    \vspace{1.0cm}

    {\large \integrantes \par}

    \vspace{1.0cm}

    {\normalsize \textbf{Asignatura:} \asignatura \par}
    {\normalsize \textbf{Docente:} \docente \par}
    {\normalsize \textbf{Fecha de presentación:} \fechaEntrega \par}

    \vspace{1.1cm}

    {\normalsize \textbf{Equipo:} \equipo \par}

    \vspace{0.9cm}

    {\normalsize \ciudad \par}

    \vspace{0.2cm}

    {\normalsize \the\year \par}
  \end{center}
\end{titlepage}


% --- Índices ----------------------------------------------------------------
\newpage
{\indicesCompactos\tableofcontents}
\newpage
{\indicesCompactos\listoftables}
\newpage
{\indicesCompactos\listoffigures}
\clearpage

% --- Cambio a numeración arábiga para el contenido principal -------------
\clearpage
\pagenumbering{arabic}

% ============================================================================
%  CONTENIDO
% ============================================================================

% --- 1. INTRODUCCIÓN --------------------------------------------------------
\caratulaCapitulo{1}{INTRODUCCI\'ON}
% ============================================================================
%  1. INTRODUCCI\'ON
%  ----------------------------------------------------------------------------
%  Contextualiza el conjunto de las actividades a4 (correo y archivos),
%  a5 (DHCP y Web), a7 (Samba 4 AD DC) y a8 (DNS del dominio y Web integrado).
% ============================================================================

\section{Introducci\'on}

En toda organizaci\'on, la infraestructura de red constituye el pilar
sobre el cual se sustentan las operaciones diarias. Un banco, con
independencia de su tama\~no, depende de que sus empleados puedan resolver
los nombres de los servicios que utilizan, de que las estaciones obtengan
autom\'aticamente su configuraci\'on de red al conectarse y de que los
sistemas internos resulten inaccesibles desde el exterior sin una
autorizaci\'on expresa. Estos requisitos, que en el uso cotidiano pasan
desapercibidos, se traducen t\'ecnicamente en un conjunto de servicios
fundamentales que deben operar de forma coordinada y segura
\cite{tanenbaum2017}.

La resoluci\'on de nombres de dominio (\textsc{dns}) constituye la pieza
central de esa coordinaci\'on. El protocolo traduce los nombres que los
usuarios recuerdan, como \texttt{web.sudoers.lan}, en las direcciones
\texttt{ip} que las m\'aquinas entienden \cite{rfc1034, rfc1035}. Sobre esa
base se articulan el resto de los servicios: el correo electr\'onico
necesita localizar al servidor de entrega a trav\'es de registros
\textsc{mx} \cite{postel1982, criden2005}; la asignaci\'on din\'amica de
direcciones evita configurar cada estaci\'n a mano \cite{rfc2131, droms1997};
y la autenticaci\'on de los usuarios se apoya en un directorio que debe ser
localizada de forma inequ\'ivoca \cite{rfc4511, rfc4120}.

\subsection{La importancia de la administraci\'on de servidores GNU/Linux}

Las organizaciones actuales operan su infraestructura sobre sistemas
libres, y no por una cuesti\'on de preferencia tecnol\'ogica sino por
razones operativas y econ\'omicas bien documentadas. Un sistema GNU/Linux
ofrece un modelo de seguridad robusto basado en la separaci\'on de
privilegios, el control de acceso discrecional y la trazabilidad de las
operaciones, principios que se aplican de forma sistem\'atica en la
administraci\'on de servidores \cite{stewart2017}. Adem\'as, la
disponibilidad del c\'odigo fuente permite auditar el comportamiento de los
servicios criticales y corregir vulnerabilidades sin depender de los
plazos del proveedor.

Administrar estos servidores exige un conjunto amplio de competencias:
configurar servicios de red, gestionar usuarios y permisos, resolver
incidencias y aplicar medidas de securizaci\'on que reduzcan la
superficie de ataque \cite{nemeth2017}. El presente trabajo final
recoge justamente ese conjunto de competencias aplicado a un escenario
empresarial concreto.

\subsection{Antecedentes y contexto del proyecto}

El proyecto se desarrollo de manera incremental a lo largo de las
actividades del curso, y cada etapa aport\'o una pieza distinta a la
arquitectura final. La primera etapa implement\'o el servicio de correo
electr\'onico con \textsc{postfix} y \textsc{dovecot}, junto con un servidor
de archivos multi-protocolo que combina \textsc{smb}, \textsc{nfs} y
\textsc{http} sobre un mismo directorio compartido. La segunda etapa
incorpor\'o el servidor de asignaci\'on din\'amica de direcciones con
\textsc{kea} y un servidor web basado en \textsc{nginx} que act\'ua como
\textit{proxy inverso} \cite{keaDocs, nginxDocs}. La tercera etapa
despleg\'o el controlador de dominio \textsc{samba 4} \textsc{active
directory}, que centraliza la identidad, la autenticaci\'on Kerberos y la
compartici\'on de archivos \cite{sambateam2024, msft_activedirectory}. La
cuarta y \'ultima etapa consolid\'o el \textsc{dns} autoritativo del dominio
servido por el propio controlador y la publicaci\'on de los servicios por
nombre a trav\'es del proxy, con cifrado \textsc{https} y una autoridad de
certificaci\'on interna \cite{iscbind9, opensslDocs}.

\subsection{Problem\'atica}

La organizaci\'on se apoya en un \'unico servidor Debian que
concentra la identidad, los archivos, el correo y la publicaci\'on web. En
ese escenario se plantearon tres carencias que limitaban la operaci\'on:

\begin{itemize}
  \item \textbf{Resoluci\'on de nombres incompleta:} sin un servidor
        \textsc{dns} propio, cada cliente deb\'ia memorizar direcciones
        \textsc{ip} o depender de archivos \texttt{/etc/hosts} dispersos,
        esquema que no escala y que impide que un mismo nombre se resuelva
        de forma consistente en toda la red.
  \item \textbf{Asignaci\'on de direcciones manual:} configurar la red de
        cada estaci\'n de forma individual consume tiempo y es propenso a
        errores, adem\'s de impedir la gesti\'on centralizada de los
        par\'ametros de la red.
  \item \textbf{Ausencia de una puerta \'unica y segura:} sin un servidor web
        que concentre la publicaci\'on de los servicios, cada aplicaci\'on
        deber\'ia exponer su propio puerto, sin una capa de cifrado
        unificada ni un punto \'unico de control y de auditor\'ia.
\end{itemize}

\subsection{Prop\'osito del trabajo final}

El prop\'osito de este trabajo es demostrar la capacidad de dise\~nar,
implementar y validar una arquitectura de red empresarial completa sobre
GNU/Linux, integrando los servicios de directorio, asignaci\'on de direcciones,
resoluci\'on de nombres, publicaci\'on web, correo electr\'onico y
archivos, con medidas de securizaci\'on aplicadas en cada uno de ellos
\cite{debianReference, hertzog2020}. El objetivo no es simplemente que los
servicios funcionen, sino que operen de forma integrada: que un cliente
reciba su configuraci\'on por \textsc{dhcp}, resuelva los nombres contra el
controlador de dominio, se autentique mediante \textsc{kerberos} y
acceda a los servicios publicados por nombre a trav\'es de un \'unico
punto de entrada cifrado.

\subsection{Marco te\'orico}

\subsubsection{El sistema de nombres de dominio}

El \textsc{dns} es un servicio distribuido y jer\'arquico que traduce
nombres de dominio a direcciones \textsc{ip} \cite{rfc1034}. En un dominio
de \textsc{active directory} no es un accesorio, sino el transporte del
directorio: los clientes lo consultan para descubrir d\'onde reside el
servidor \textsc{ldap}, el \textsc{kerberos} o el propio controlador
mediante registros \textsc{srv} \cite{rfc2782}. Por esta raz\'on, un
controlador de dominio debe ser autoritativo de su zona
\cite{msft_activedirectory}. Los registros empleados en el proyecto son
los de tipo \textsc{a}, que asocian un nombre a una direcci\'on; los
\textsc{srv}, que indican qu\'e servidor atiende qu\'e servicio en qu\'e
puerto; y los registros \textsc{ns} y \textsc{soa}, que declaran la
autoridad y la vigencia de la zona.

\subsubsection{La asignaci\'on din\'amica de direcciones}

El protocolo \textsc{dhcp} permite que un servidor asigne direcciones
\textsc{ip} de forma autom\'atica a los clientes que se conectan
\cite{rfc2131}. El proceso sigue una secuencia de cuatro mensajes en la que
el cliente descubre, ofrece, solicita y confirma una direcci\'on. Adem\'as
de la direcci\'on, el servidor entrega par\'ametros de configuraci\'on
como la m\'ascara de subred, la puerta de enlace y el servidor
\textsc{dns}, lo que evita la intervenci\'on manual \cite{droms1997}. El
servidor \textsc{kea} se escogi\'o por su soporte activo y su
configuraci\'on basada en \textsc{json} \cite{keaDocs}.

\subsubsection{El servidor web y el proxy inverso}

Un servidor web atiende peticiones \textsc{http} y entrega documentos a
los navegadores. En el proyecto se emplea \textsc{nginx} con una doble
funci\'on: servir el contenido del portal y actuar como \textit{proxy
inverso}, es decir, recibir las peticiones y reenviarlas al servicio
interno correspondiente seg\'un el nombre solicitado \cite{nginxDocs}. El
proxy termina adem\'as la conexi\'on \textsc{tls}, de modo que el
cifrado se aplica en el punto \'unico de entrada mientras que las
comunicaciones internas se realizan en claro.

\subsubsection{La identidad centralizada}

El \textsc{active directory} centraliza la gesti\'on de usuarios, grupos y
permisos sobre la base de \textsc{ldap} \cite{rfc4511}, y delega la
autenticaci\'on en \textsc{kerberos}, cuyo protocolo de tickets evita
transmitir la clave del usuario en cada petici\'on \cite{rfc4120}. Esta
arquitectura es la que permite, en lugar de credenciales dispersas por
cada servicio, una gesti\'on \'unica y coherente de la seguridad
\cite{sambateam2024}.

\subsection{Estructura del documento}

El presente informe se organiza en seis cap\'itulos. El cap\'itulo
siguiente enuncia los objetivos del trabajo y su alcance. El tercero
documenta el diagrama arquitect\'onico final, con la segmentaci\'on de la
red, los roles de cada servidor y la justificaci\'on del dise\~no. El
cuarto desarrolla t\'ecnicamente cada servicio desplegado, siguiendo el
orden en que se implementaron. El quinto expone los desaf\'ios
encontrados durante la ejecuci\'on y las soluciones aplicadas. El sexto
expone las conclusiones, los aprendizajes y las recomendaciones para
futuras implementaciones.


% --- 2. OBJETIVOS DEL PROYECTO ----------------------------------------------
\caratulaCapitulo{2}{OBJETIVOS DEL PROYECTO}
% ============================================================================
%  2. OBJETIVOS DEL PROYECTO
% ============================================================================

\section{Objetivos del proyecto}

\subsection{Objetivo general}

Diseñar e implementar una arquitectura empresarial GNU/Linux funcional y
segura, integrando servicios de directorio (Active Directory), DHCP, DNS,
Web, correo electrónico y archivos, con medidas de securización aplicadas en
cada uno de los servicios desplegados \cite{nemeth2017}.

\subsection{Objetivos específicos}

\begin{enumerate}
  \item Diseñar una arquitectura empresarial funcional y segura.
  \item Implementar un servidor de correo electrónico (Postfix/Dovecot) y un
        servidor de archivos multi-protocolo (Samba/NFS/HTTP) con resolución
        de nombres asociada \cite{postel1982}.
  \item Configurar un servidor DHCP (Kea) y un servidor Web con proxy
        inverso y soporte TLS \cite{rfc2131, keaDocs}.
  \item Desplegar un controlador de dominio Samba 4 Active Directory con
        DNS, DHCP y autenticación Kerberos \cite{sambateam2024, rfc4120}.
  \item Validar el funcionamiento integral de la infraestructura mediante
        pruebas funcionales desde los clientes de la red.
\end{enumerate}


% --- 3. DIAGRAMA ARQUITECTÓNICO FINAL ---------------------------------------
\caratulaCapitulo{3}{DIAGRAMA ARQUITECT\'ONICO FINAL}
% ============================================================================
%  3. DIAGRAMA ARQUITECTÓNICO FINAL
%  ----------------------------------------------------------------------------
%  Fuente: TSO-II-final/docs/a10.md, ajustado con los datos de red
%  vigentes del proyecto (AGENTS.md).
% ============================================================================

\section{Diagrama arquitectónico final}

La infraestructura del banco Sudoers se diseñó sobre una red local única
(\texttt{192.168.0.0/24}) que se segmenta l\'ogicamente por roles y nivel de
acceso, de modo que cada zona expone \'unicamente los servicios que
necesita. Esta secci\'on documenta la segmentaci\'on adoptada, los roles
asignados a cada servidor, el flujo de comunicaci\'on resultante y la
justificaci\'on t\'ecnica del dise\~no.

\subsection{Segmentaci\'on de red}

La red se divide en cuatro zonas funcionales, delimitadas por rangos de
direcciones y pol\'iticas de filtrado distintas, tal como se detalla en la
tabla \ref{tab:zonas-red}.

\begin{table}[H]
\centering
\caption{Segmentaci\'on de la red local por zonas funcionales}
\label{tab:zonas-red}
\contenidoConFuente{%
  \begin{tablaAPA}
    \begin{tabular}{@{}p{2.4cm}p{4.2cm}p{3.4cm}p{4cm}@{}}
    \toprule
    \textbf{Zona} & \textbf{Rango} & \textbf{Funci\'on} & \textbf{Equipos} \\
    \midrule
    A & \texttt{192.168.0.1} & Per\'imeter y salida a Internet & Router TP-Link TL-WR850N \\
    B & \texttt{192.168.0.2} & Host servidor de dominio y servicios & \texttt{dc1} (Debian 12) \\
    B & \texttt{192.168.0.3} & Monitoreo de infraestructura & \texttt{zabbix} (Debian 13) \\
    C & \texttt{.60}--\texttt{.99} & Puestos operativos y cajas & Terminales Linux Kiosk \\
    D & \texttt{.50}--\texttt{.52} & Administraci\'on de TI & \texttt{pc-joaquin}, \texttt{pc-david}, \texttt{pc-nicolas} \\
    D & \texttt{.100}--\texttt{.199} & Administraci\'on y gerencia (din\'amico) & Estaciones Windows 10/11 \\
    \bottomrule
    \end{tabular}
  \end{tablaAPA}%
}{Elaboraci\'on propia.}
\end{table}

El servicio DHCP del router (Zona A) se desactiva de forma deliberada para
que el servidor Kea sea el \'unico emisor de ofertas de la red, evitando as\'\i
conflicto de doble respuesta \textsc{Discover} \cite{rfc2131}.

\subsection{Servidores y roles}

La tabla \ref{tab:servidores} detalla cada servidor, su direcci\'on y el rol
que ejerce dentro de la arquitectura.

\begin{table}[H]
\centering
\caption{Servidores, direcciones y roles en la arquitectura final}
\label{tab:servidores}
\contenidoConFuente{%
  \begin{tablaAPA}
    \begin{tabular}{@{}p{2.2cm}p{3.2cm}p{2.4cm}p{5.8cm}@{}}
    \toprule
    \textbf{Host} & \textbf{IP} & \textbf{SO} & \textbf{Rol} \\
    \midrule
    \texttt{dc1} & \texttt{192.168.0.2} & Debian 12 & Controlador de dominio Samba 4, DNS, NTP, firewall, servicios en Docker \\
    \texttt{zabbix} & \texttt{192.168.0.3} & Debian 13 & Servidor de monitoreo Zabbix (server, web y PostgreSQL) \\
    Router & \texttt{192.168.0.1} & Firmware & Puerta de enlace, NAT a Internet y segmentaci\'on \\
    \bottomrule
    \end{tabular}
  \end{tablaAPA}%
}{Elaboraci\'on propia.}
\end{table}

El host \texttt{dc1} concentra una arquitectura de tres capas de
procesamiento. La capa nativa ejecuta el controlador de dominio Samba 4, que
integra en un \'unico daemon los servicios de LDAP, Kerberos, DNS y SMB
\cite{sambateam2024}, junto con \texttt{chrony} para la sincronizaci\'on
horaria --imprescindible para Kerberos, que rechaza los tickets con un
desfase de reloj superior a cinco minutos--, el firewall \texttt{nftables} y el servicio
\texttt{sshd} endurecido. La segunda capa corresponde a los contenedores
Docker, donde residen el proxy inverso Nginx, el servidor DHCP Kea en modo
de red \texttt{host}, la base de datos PostgreSQL, el servicio de correo
Postfix y Dovecot, el gestor de impresi\'on CUPS y el panel de gesti\'on
Portainer. La tercera capa constituye el servidor \texttt{zabbix}, que
opera como m\'aquina independiente y es alcanzable \'unicamente a trav\'es
del proxy inverso.

La decisi\'on de ejecutar el controlador de dominio de forma nativa, y no
contenedorizado, responde a la fragilidad de Samba frente al aislamiento de
redes y sistemas de archivos en contenedores \cite{sambateam2024}.

\subsection{Flujo de comunicaci\'on entre servicios}

La figura \ref{fig:arquitectura-zonas} representa la topolog\'ia general de
la red y las relaciones entre el per\'imeter, el host de servicios, el
servidor de monitoreo y las zonas de puestos de trabajo.

\capturaAPA{fig:arquitectura-zonas}{Diagrama de la arquitectura de red segmentada por zonas y roles de acceso}{captura-arquitectura-red.jpeg}{Diagrama de la topolog\'ia general de la red del banco Sudoers, con la segmentaci\'on en zonas A (per\'imeter), B (servidores \texttt{dc1} y \texttt{zabbix}), C (puestos de cajas) y D (administraci\'on y gerencia), junto con los flujos de comunicaci\'on entre ellas.}

La matriz de la tabla \ref{tab:matriz-acceso} formaliza el flujo anterior:
para cada servicio se indican el puerto y protocolo expuestos, la capa de
ejecici\'on y el nivel de acceso permitido.

\begin{table}[H]
\centering
\caption{Matriz de acceso y visibilidad de los servicios}
\label{tab:matriz-acceso}
\contenidoConFuente{%
  \begin{tablaAPA}
    \begin{tabular}{@{}p{4.9cm}p{2.4cm}p{2.1cm}p{4.5cm}@{}}
    \toprule
    \textbf{Servicio} & \textbf{Puerto} & \textbf{Capa} & \textbf{Acceso} \\
    \midrule
    Samba DNS / LDAP / Kerberos & \texttt{53}, \texttt{88}, \texttt{389} & Nativa & Todo el dominio (validaci\'on) \\
    Recursos compartidos SMB & \texttt{445/tcp} & Nativa & Filtrado por grupos de usuario AD \\
    SSH / nftables & \texttt{22/tcp} & Nativa & Exclusivo de administradores TI (\texttt{.50}--\texttt{.52}) \\
    Nginx proxy inverso & \texttt{80}, \texttt{443/tcp} & Docker & LAN operativa y gerencia \\
    Postfix / Dovecot & \texttt{25}, \texttt{143}, \texttt{587} & Docker & Usuarios autenticados del dominio \\
    Kea DHCP & \texttt{67/udp} & Docker (host) & Difusi\'on de toda la LAN \\
    PostgreSQL & \texttt{5432/tcp} & Docker & Aislado (solo Core Bancario) \\
    CUPS & \texttt{631/tcp} & Docker & Puestos de cajas y administraci\'on \\
    Portainer Web & \texttt{9443/tcp} & Docker & Exclusivo de administradores TI \\
    Zabbix (agente pasivo) & \texttt{10050/tcp} & Servidor \texttt{zabbix} & Clientes y host \\
    \bottomrule
    \end{tabular}
  \end{tablaAPA}%
}{Elaboraci\'on propia.}
\end{table}

\subsection{Justificaci\'on t\'ecnica del dise\~no}

El dise\~no se apoya en cinco decisiones t\'ecnicas fundamentadas.

\begin{enumerate}
  \item \textbf{Segmentaci\'on por roles.} Dividir la red en zonas con
        pol\'iticas de filtrado distintas reduce la superficie de ataque:
        los puestos de cajas no requieren acceso al panel de administraci\'on
        de contenedores ni a la consola del servidor de monitoreo.
  \item \textbf{Seguridad por defecto.} La pol\'itica \textsc{drop} por
        defecto de \texttt{nftables} y la restricci\'n del puerto
        \texttt{22/tcp} \'unicamente a los rangos de administraci\'on
        implementan el principio de m\'inimo privilegio
        \cite{stewart2017}.
  \item \textbf{Identidad centralizada.} Delegar la autenticaci\'on en
        Samba 4 AD DC sobre Kerberos y LDAP permite una gesti\'on \'unica de
        usuarios, grupos y permisos, en lugar de credenciales dispersas por
        cada servicio \cite{rfc4120, rfc4511}.
  \item \textbf{Aislamiento de datos sensibles.} PostgreSQL queda
        confinado a la red interna de Docker y el servicio de monitoreo se
        publica \'unicamente a trav\'es del proxy inverso, de modo que
        ninguno es alcanzable de forma directa desde la LAN.
  \item \textbf{Un \'unico emisor de direcciones.} Desactivar el DHCP del
        router y centralizar la asignaci\'on en Kea elimina la ambig\"uedad
        en la obtenci\'on de la direcci\'on, puerta de enlace y servidor DNS
        por parte de los clientes \cite{rfc2131, keaDocs}.
\end{enumerate}

Estas decisiones culminan, en conjunto, en una arquitectura que puede escalar
con nuevos servicios sobre la misma base de segmentaci\'on y controles de
acceso, principios que se verifican de forma pr\'actica en el cap\'itulo de
desarrollo t\'ecnico.


% --- 4. DESARROLLO TÉCNICO DEL PROYECTO -------------------------------------
\caratulaCapitulo{4}{DESARROLLO T\'ECNICO DEL PROYECTO}
% ============================================================================
%  4. DESARROLLO T\'ECNICO DEL PROYECTO
%  ----------------------------------------------------------------------------
%  Integraci\'on del contenido de las actividades a4 (correo y archivos),
%  a5 (DHCP y Web) y a7 (Samba 4 AD DC).
% ============================================================================

\section{Desarrollo t\'ecnico del proyecto}

Este cap\'itulo documenta la implementaci\'on de los servicios de
infraestructura del banco Sudoers. El desarrollo se organiz\'o en tres
bloques correspondientes a las actividades del curso: primero el servicio de
correo electr\'onico y el de archivos multi-protocolo, luego el servidor DHCP
y el servidor Web con proxy inverso, y finalmente el despliegue del
controlador de dominio Samba 4 Active Directory que integra la identidad de
toda la red.

Cada bloque describe desde la instalaci\'on de los paquetes necesarios hasta
las pruebas funcionales ejecutadas desde los clientes de la LAN. Las
direcciones IP citadas corresponden al plan de direccionamiento final
documentado en el cap\'itulo de arquitectura: el host \texttt{dc1} opera en
\texttt{192.168.0.2}, el servidor de monitoreo \texttt{zabbix} en
\texttt{192.168.0.3} y el router en \texttt{192.168.0.1}.

% ============================================================================
%  BLOQUE 1: SERVIDOR DE CORREO ELECTR\'ONICO
% ============================================================================
\section{Servidor de correo electr\'onico}

El servicio de correo se implement\'o sobre Debian 12 mediante dos
componentes instalados de forma nativa: Postfix como agente de
transferencia de correo (\textsc{mta}) y Dovecot como agente de acceso a
buzones (\textsc{mda}). La protecci\'on de las conexiones se realiza con
\textsc{starttls} y certificados auto-firmados \cite{postel1982}.

\subsection{Arquitectura del servicio}

Postfix gestiona el env\'io y la recepci\'on mediante SMTP, habilitando el
puerto 25 para la transferencia entre servidores y el 587 para el env\'io
autenticado desde clientes. Dovecot expone los buzones mediante IMAP
(puerto 143) y POP3 (puerto 110), y valida las credenciales de los usuarios
a trav\'es del subsistema SASL.

La autenticaci\'on se delega de Postfix a Dovecot mediante un socket Unix
(\texttt{/var/run/dovecot/private/auth}). El cliente remite sus
credenciales en el puerto 587, Postfix consulta al socket de Dovecot y s\'olo
tras la validaci\'on exitosa acepta el mensaje. La figura
\ref{fig:arquitectura-correo} muestra el esquema del servicio.

\figuraAPA{fig:arquitectura-correo}{Arquitectura del servicio de correo: Postfix, Dovecot y DNS}{%
  \begin{tikzpicture}[
      node distance=1.2cm and 1.8cm,
      servicio/.style={draw, rounded corners, minimum width=2.6cm,
                       minimum height=0.9cm, align=center, font=\small},
      cliente/.style={draw, circle, minimum size=1cm, font=\small},
      linea/.style={-Stealth, thick}
    ]
    \node[servicio] (postfix) {Postfix\\{\footnotesize SMTP}};
    \node[servicio, right=of postfix] (dovecot) {Dovecot\\{\footnotesize IMAP/POP3}};
    \node[cliente, left=of postfix] (c1) {C1};
    \node[cliente, below=of c1] (c2) {C2};
    \node[servicio, below=of postfix, yshift=-0.4cm] (dns) {BIND9\\{\footnotesize DNS}};
    \draw[linea] (c1.east) -- node[above,font=\scriptsize]{587} (postfix.west);
    \draw[linea] (c2.east) -| (postfix.west);
    \draw[linea] (postfix.east) -- node[above,font=\scriptsize]{SASL} (dovecot.west);
    \draw[linea, dashed] (postfix.south) -- (dns.north);
  \end{tikzpicture}%
}{Elaboraci\'on propia.}

\subsection{Instalaci\'on de paquetes}

\begin{lstlisting}[style=terminal]
$ apt update
$ apt install -y postfix dovecot-core dovecot-lmtpd \
  libsasl2-modules openssl
\end{lstlisting}

El paquete \texttt{postfix} aporta el daemon \texttt{master} y el gestor
\texttt{postconf}; \texttt{dovecot-core} y \texttt{dovecot-lmtpd} instalan el
servidor de buzones con soporte \textsc{lmtp}; \texttt{libsasl2-modules}
habilita los mecanismos de autenticaci\'on; y \texttt{openssl} se emplea
para emitir los certificados \textsc{tls} \cite{opensslDocs}.

\subsection{Configuraci\'on del servidor DNS}

Antes de configurar el correo se estableci\'o la resoluci\'on de nombres del
dominio \texttt{sudoers.lan}, con un registro \textsc{mx} que apunta al
servidor de correo y su registro \textsc{a} correspondiente. La
verificaci\'on se realiz\'o con la utilidad \texttt{dig}:

\begin{lstlisting}[style=terminal]
$ dig +short MX sudoers.lan
10 mail.sudoers.lan.

$ dig +short mail.sudoers.lan
192.168.0.2
\end{lstlisting}

\capturaAPA{fig:dig-mx}{Verificaci\'on del registro \textsc{mx} del dominio \texttt{sudoers.lan} con \texttt{dig}}{captura-dig-mx.png}{Captura de pantalla de la salida de \texttt{dig +short MX sudoers.lan} y \texttt{dig +short mail.sudoers.lan} en la terminal, mostrando el registro MX y la direcci\'on IP del servidor de correo.}

\subsection{Configuraci\'on del servidor SMTP}

La configuraci\'on de Postfix reside en \texttt{/etc/postfix/main.cf}. Los
par\'ametros relevantes se resumen en la tabla \ref{tab:parametros-postfix}.

\begin{table}[H]
\centering
\caption{Par\'ametros principales de Postfix}
\label{tab:parametros-postfix}
\contenidoConFuente{%
  \begin{tablaAPA}
    \begin{tabular}{@{}p{5.8cm}p{8cm}@{}}
    \toprule
    \textbf{Par\'ametro} & \textbf{Valor} \\
    \midrule
    \texttt{myhostname} & \texttt{mail.sudoers.lan} \\
    \texttt{mydomain} & \texttt{sudoers.lan} \\
    \texttt{mynetworks} & \texttt{127.0.0.0/8, 10.0.0.0/8,}\newline
                          \texttt{172.16.0.0/12, 192.168.0.0/16} \\
    \texttt{home\_mailbox} & \texttt{Maildir/} \\
    \texttt{smtpd\_sasl\_type} & \texttt{dovecot} \\
    \texttt{smtpd\_sasl\_path} & \texttt{private/auth} \\
    \texttt{smtpd\_sasl\_auth\_enable} & \texttt{yes} \\
    \texttt{smtpd\_sasl\_security\_options} & \texttt{noanonymous} \\
    \texttt{smtpd\_use\_tls} & \texttt{yes} \\
    \texttt{smtpd\_tls\_security\_level} & \texttt{may} \\
    \bottomrule
    \end{tabular}
  \end{tablaAPA}%
}{Elaboraci\'on propia.}
\end{table}

El almacenamiento en formato \texttt{Maildir/} garantiza que cada mensaje
se guarde como un archivo independiente, lo que evita bloqueos al escribir
y simplifica el respaldo. La opci\'on \texttt{noanonymous} impide el env\'io
sin autenticaci\'on.

\capturaAPA{fig:postconf}{Par\'ametros efectivos de Postfix con \texttt{postconf -n}}{captura-postconf.png}{Captura de pantalla de la terminal ejecutando \texttt{postconf -n}, mostrando la configuraci\'on activa del servidor de correo, incluidos los par\'ametros de identidad, redes y SASL.}

\subsection{Configuraci\'on del servidor IMAP/POP3}

Dovecot se configura mediante los archivos del directorio
\texttt{/etc/dovecot/conf.d/}. La configuraci\'on resultante se verificó con
el comando \texttt{doveconf -n}.

\capturaAPA{fig:doveconf}{Verificaci\'on de la configuraci\'on de Dovecot con \texttt{doveconf -n}}{captura-doveconf.png}{Captura de pantalla de la terminal ejecutando \texttt{doveconf -n}, con los par\'ametros activos de Dovecot referidos a protocolos, ubicaci\'on del buz\'on y cifrado.}

\subsection{Pruebas de conectividad y de env\'io}

Desde el equipo cliente se verific\'o la apertura de los puertos del
servicio de correo:

\begin{lstlisting}[style=terminal]
$ for p in 25 110 143 587 993 995; do
    timeout 3 bash -c "echo > /dev/tcp/192.168.0.2/$p" \
      && echo "puerto $p: ABIERTO"
  done
puerto 25: ABIERTO
puerto 587: ABIERTO
\end{lstlisting}

\capturaAPA{fig:puertos-abiertos}{Comprobaci\'on de conectividad a los puertos del servidor de correo}{captura-puertos.png}{Captura de pantalla mostrando la salida del bucle que verifica los puertos 25, 110, 143, 587, 993 y 995 desde la m\'aquina cliente.}

La negociaci\'on \textsc{starttls} se comprob\'o con \texttt{openssl},
que muestra el sujeto del certificado y el c\'odigo de verificaci\'on 18
correspondiente a un certificado auto-firmado:

\begin{lstlisting}[style=terminal]
$ openssl s_client -connect 192.168.0.2:587 -starttls smtp
subject=CN=mail.sudoers.lan, O=TSO-II, C=AR
Verify return code: 18 (self-signed certificate)
\end{lstlisting}

\capturaAPA{fig:starttls}{Negociaci\'on \textsc{starttls} con el servidor SMTP}{captura-starttls.png}{Captura de pantalla de \texttt{openssl s\_client} mostrando el sujeto del certificado y el c\'odigo de verificaci\'on 18, propio de un certificado auto-firmado.}

El env\'io autenticado se valid\'o con \texttt{swaks}, que completa el
intercambio \textsc{smtp} y obtiene la respuesta \texttt{250} de aceptaci\'on
del mensaje:

\begin{lstlisting}[style=terminal]
$ swaks --to joaquin@sudoers.lan \
    --from nicolas@sudoers.lan \
    --server 192.168.0.2 --port 587 \
    --auth LOGIN --auth-user nicolas \
    --header "Subject: Prueba" \
    --body "Hola Joaquin, esto es una prueba."
\end{lstlisting}

\capturaAPA{fig:swaks}{Env\'io de correo autenticado con \texttt{swaks}}{captura-swaks.png}{Captura de pantalla de \texttt{swaks} mostrando el intercambio SMTP, la autenticaci\'on exitosa y el c\'odigo 250 de aceptaci\'on del mensaje.}

La recepci\'on se confirm\'o inspeccionando el buz\'on del destinatario,
donde el mensaje aparece como un archivo individual dentro del directorio
\texttt{new/} del \texttt{Maildir}:

\begin{lstlisting}[style=terminal]
$ ls /home/joaquin/Maildir/new/
1627...joaquin@sudoers.lan
\end{lstlisting}

\capturaAPA{fig:maildir-receptor}{Buz\'on \texttt{Maildir} del destinatario con el mensaje recibido}{captura-maildir.png}{Captura de pantalla listando el contenido de \texttt{/home/joaquin/Maildir/new/}, donde aparece el archivo del correo recibido con el nombre generado por Postfix.}

\subsection{Acceso desde clientes}

El acceso desde un cliente gr\'afico como Thunderbird se realiz\'o con la
configuraci\'on que se detalla en la tabla \ref{tab:config-thunderbird}.

\begin{table}[H]
\centering
\caption{Configuraci\'on del cliente de correo Thunderbird}
\label{tab:config-thunderbird}
\contenidoConFuente{%
  \begin{tablaAPA}
    \begin{tabular}{@{}p{5.4cm}p{8.4cm}@{}}
    \toprule
    \textbf{Par\'ametro} & \textbf{Valor} \\
    \midrule
    Servidor de entrada (IMAP) & \texttt{mail.sudoers.lan} \\
    Puerto de entrada & \texttt{143} \\
    Seguridad de entrada & \textsc{starttls} \\
    Servidor de salida (SMTP) & \texttt{mail.sudoers.lan} \\
    Puerto de salida & \texttt{587} \\
    Seguridad de salida & \textsc{starttls} \\
    Autenticaci\'on & Contrase\~na normal \\
    Usuario & \texttt{joaquin} \\
    \bottomrule
    \end{tabular}
  \end{tablaAPA}%
}{Elaboraci\'on propia.}
\end{table}

\capturaAPA{fig:thunderbird}{Cliente Thunderbird conectado al servidor IMAP con los mensajes recibidos}{captura-thunderbird.png}{Captura de pantalla del cliente Thunderbird mostrando la cuenta \texttt{joaquin@sudoers.lan} conectada y la bandeja de entrada con el mensaje de prueba.}

% ============================================================================
%  BLOQUE 2: SERVIDOR DE ARCHIVOS MULTI-PROTOCOLO
% ============================================================================
\section{Servidor de archivos multi-protocolo}

El servicio de archivos integra tres protocolos --\textsc{smb/cifs} con
Samba, \textsc{nfs} con el m\'odulo del n\'ucleo y \textsc{http} con
Filebrowser-- que operan sobre un mismo directorio compartido
(\texttt{/srv/share}). De este modo, un archivo creado desde un protocolo es
inmediatamente accesible desde los dem\'as, sin duplicar la informaci\'on
almacenada.

\subsection{Configuraci\'on de Samba}

La configuraci\'on reside en \texttt{/etc/smb.conf}. Las directivas
esenciales y su justificaci\'on se resumen en la tabla
\ref{tab:directivas-samba}.

\begin{table}[H]
\centering
\caption{Directivas principales de la configuraci\'on de Samba}
\label{tab:directivas-samba}
\contenidoConFuente{%
  \begin{tablaAPA}
    \begin{tabular}{@{}p{5.4cm}p{8.4cm}@{}}
    \toprule
    \textbf{Directiva} & \textbf{Justificaci\'on} \\
    \midrule
    \texttt{security = user} & Exige autenticaci\'on con usuario y clave registrados \\
    \texttt{map to guest = bad user} & Rechaza usuarios inv\'alidos sin otorgar acceso de invitado \\
    \texttt{guest ok = no} & Proh\'ibe el acceso an\'onimo al recurso \\
    \texttt{valid users} & Lista blanca de los usuarios autorizados \\
    \texttt{create mask / directory mask} & \texttt{0666} y \texttt{0777}: los archivos de un protocolo quedan accesibles a los dem\'as \\
    \bottomrule
    \end{tabular}
  \end{tablaAPA}%
}{Elaboraci\'on propia.}
\end{table}

\begin{lstlisting}[style=terminal]
[compartido]
    path = /srv/share
    browseable = yes
    writable = yes
    guest ok = no
    valid users = nicolas joaquin david
    create mask = 0666
    directory mask = 0777
\end{lstlisting}

El acceso desde la terminal se verific\'o con \texttt{smbclient}, que
autentica al usuario y permite listar y crear contenido en el recurso
compartido:

\capturaAPA{fig:archivos-smbclient}{Conexi\'on al recurso compartido con \texttt{smbclient}}{captura-archivos-smbclient.png}{Captura de pantalla de una sesi\'on \texttt{smbclient} contra el recurso compartido, autenticada con el usuario david y mostrando la lista de archivos y la creaci\'on de un directorio.}

\subsection{Configuraci\'on de NFS}

La exportaci\'n se define en \texttt{/etc/exports}. La opci\'on
\texttt{fsid=0} establece el directorio como ra\'iz virtual del espacio
\textsc{nfsv4}, por lo que el cliente monta la ruta remota \texttt{/} y no
\texttt{/srv/share}:

\begin{lstlisting}[style=terminal]
/srv/share *(rw,sync,no_subtree_check,no_root_squash,fsid=0)

$ showmount -e 192.168.0.2
$ sudo mount -t nfs4 192.168.0.2:/ /mnt/servidor_archivos
\end{lstlisting}

La opci\'on \texttt{sync} prioriza la integridad de los datos al confirmar
en disco antes de responder, mientras que \texttt{no\_subtree\_check}
optimiza el rendimiento en las renombraciones.

\capturaAPA{fig:archivos-nfs}{Verificaci\'on de exportaciones y montaje \textsc{nfs}}{captura-archivos-nfs.png}{Captura de pantalla mostrando la salida de \texttt{showmount -e} y el montaje de \texttt{/} en el punto de acceso del cliente.}

\subsection{Interfaz web de archivos}

Filebrowser expone el mismo directorio \texttt{/srv/share} a trav\'es de
una interfaz web accesible desde cualquier navegador, sin instalar software
adicional en el cliente. La inicializaci\'on de su base de datos y la
creaci\'on de usuarios se realizan al desplegar el servicio:

\begin{lstlisting}[style=terminal]
$ filebrowser config init --database /database/filebrowser.db
$ filebrowser config set --database /database/filebrowser.db \
    --root /srv/share --address 0.0.0.0 --port 80 \
    --auth.method json
\end{lstlisting}

\capturaAPA{fig:archivos-filebrowser}{Interfaz web de Filebrowser autenticada}{captura-archivos-filebrowser.png}{Captura de pantalla del navegador mostrando la pantalla de inicio de sesi\'on y el explorador de archivos sobre el directorio \texttt{/srv/share}.}

% ============================================================================
%  BLOQUE 3: SERVIDOR DHCP
% ============================================================================
\section{Servidor \textsc{dhcp}}

El servicio de asignaci\'on din\'amica de direcciones se implement\'o con
Kea, el servidor \textsc{dhcp} de ISC y sucesor del cl\'asico
\texttt{dhcpd}, seleccionado por su soporte activo, su configuraci\'on
basada en \textsc{json} y su manejo eficiente de las concesiones
\cite{rfc2131, keaDocs}.

\subsection{Arquitectura y segmentaci\'on de direcciones}

El servicio se despleg\'o de manera centralizada sobre el host \texttt{dc1}.
Para que \'unicamente este servidor otorgue direcciones, se desactiv\'o el
servicio \textsc{dhcp} del router, evitando que dos servidores respondieran a
las peticiones \textsc{discover} con ofertas distintas.

La red se dividi\'o en dos segmentos l\'ogicos seg\'un la pol\'itica del
proyecto: las direcciones inferiores a \texttt{192.168.0.110} corresponden a
los equipos de rol, que requieren \textsc{ip} fija, mientras que el rango
superior se reserva para la asignaci\'on din\'amica de los equipos de
usuario. La figura \ref{fig:arquitectura-dhcp} representa el esquema del
servicio.

\figuraAPA{fig:arquitectura-dhcp}{Arquitectura del servicio \textsc{dhcp}: clientes, router y servidor}{%
  \begin{tikzpicture}[
      node distance=1.2cm and 1.6cm,
      ser/.style={draw, rounded corners, minimum width=2.6cm,
                  minimum height=0.9cm, align=center, font=\small},
      cli/.style={draw, minimum width=2.4cm, minimum height=0.8cm,
                  align=center, font=\small},
      flujo/.style={-Stealth, thick}
    ]
    \node[ser] (kea) {Kea DHCP\\\texttt{192.168.0.2}};
    \node[cli, left=of kea, yshift=0.7cm] (c1) {Equipo de usuario\\\texttt{.100+}};
    \node[cli, left=of kea, yshift=-0.7cm] (c2) {Equipo de rol\\\texttt{.50--.99}};
    \node[ser, right=of kea] (router) {Router\\\texttt{192.168.0.1}\\\footnotesize DHCP off};
    \node[ser, below=0.9cm of kea] (dns) {DNS\\\texttt{192.168.0.2}\\\footnotesize mismo host};
    \draw[flujo] (c1.east) -- (kea.west);
    \draw[flujo] (c2.east) -- (kea.west);
    \draw[flujo] (kea.east) -- (router.west);
    \draw[flujo, dashed] (dns.north) -- (kea.south);
  \end{tikzpicture}%
}{Elaboraci\'on propia.}

\subsection{Instalaci\'on y configuraci\'on del servidor}

\begin{lstlisting}[style=terminal]
$ apt update
$ apt install -y kea-dhcp4-server
\end{lstlisting}

La configuraci\'on reside en \texttt{/etc/kea/kea-dhcp4.conf}, en formato
\textsc{json}. El bloque \texttt{subnet4} define el \'ambito, el conjunto de
direcciones y las opciones de red que se entregan al cliente:

\begin{lstlisting}[style=terminal]
"interfaces-config": { "interfaces": ["*"] },
"lease-database": { "type": "memfile", "lfc-interval": 3600 },
"subnet4": [{
  "subnet": "192.168.0.0/24",
  "pools": [{ "pool": "192.168.0.100 - 192.168.0.199" }],
  "option-data": [
    { "name": "domain-name-servers", "data": "192.168.0.2" },
    { "name": "domain-name", "data": "sudoers.lan" },
    { "name": "routers", "data": "192.168.0.1" }
  ]
}]
\end{lstlisting}

El par\'ametro \texttt{lfc-interval} indica cada cu\'antos segundos se
compacta el archivo de concesiones. Las opciones de \texttt{option-data}
garantizan que cada cliente reciba, sin intervenci\'on manual, la m\'ascara
de subred, la puerta de enlace, el servidor \textsc{dns} y el dominio de
b\'usqueda del proyecto.

\capturaAPA{fig:dhcp-conf}{Archivo de configuraci\'on \texttt{kea-dhcp4.conf} con el \'ambito y las opciones de red}{captura-dhcp-conf.png}{Captura de pantalla del archivo de configuraci\'on de Kea mostrando el bloque \texttt{subnet4} con el \'ambito din\'amico, el servidor DNS, el dominio y la puerta de enlace definidos.}

\subsection{Desactivaci\'on del \textsc{dhcp} del router}

El router de la red manten\'ia su propio servicio \textsc{dhcp} activo,
lo que generaba un conflicto: tanto el router como Kea respond\'ian a los
paquetes \textsc{discover}. Para garantizar que el \'unico servidor
\textsc{dhcp} de la red sea el implementado en el proyecto, se accedi\'o a
la interfaz de administraci\'on del router y se deshabilit\'o el servicio en
la \textsc{lan}, manteniendo la configuraci\'on de la puerta de enlace.

\capturaAPA{fig:router-disable}{Interfaz de administraci\'on del router con el servidor \textsc{dhcp} deshabilitado}{captura-router-disable.png}{Captura de pantalla de la configuraci\'on \textsc{lan} del router TP-Link mostrando la opci\'on \textsc{dhcp server} en estado deshabilitado.}

\subsection{Verificaci\'on y pruebas de asignaci\'on}

Antes de probar con clientes reales se valid\'o la sintaxis de la
configuraci\'on y el estado del servicio:

\begin{lstlisting}[style=terminal]
$ kea-dhcp4 -t /etc/kea/kea-dhcp4.conf
Configuration check successful
$ systemctl restart kea-dhcp4-server
$ systemctl status kea-dhcp4-server
\end{lstlisting}

\capturaAPA{fig:estado-kea}{Estado del servicio \texttt{kea-dhcp4-server} en el servidor}{captura-systemctl-kea.png}{Captura de pantalla de la terminal mostrando el estado activo del servicio \texttt{kea-dhcp4-server} con la ruta de la configuraci\'on cargada.}

Desde un equipo cliente se solicit\'o una concesion y se verific\'o la
configuraci\'on recibida:

\begin{lstlisting}[style=terminal]
$ sudo dhclient -v eth0
$ ip a
$ ip route
\end{lstlisting}

\capturaAPA{fig:cliente-ip}{Direcci\'on asignada al cliente tras el arranque de \texttt{dhclient}}{captura-cliente-ip.png}{Captura de pantalla de la salida de \texttt{ip a} en el cliente mostrando la direcci\'on IP dentro del \'ambito din\'amico, la m\'ascara \texttt{/24} y el estado de la interfaz.}

Cada asignaci\'n queda registrada en el archivo de concesiones, que
identifica la direcci\'n entregada, la direcci\'n \textsc{mac} del cliente y
los tiempos de validez:

\capturaAPA{fig:leases}{Concesiones registradas en \texttt{kea-leases4.csv}}{captura-leases.png}{Captura de pantalla del archivo de concessions mostrando las concesiones activas: direcci\'on IP, direcci\'on MAC del cliente, \texttt{hostname} y tiempos de renovaci\'on y expiraci\'on.}

Finalmente se verific\'o la conectividad del cliente hacia la puerta de
enlace y hacia el servidor \textsc{dns}, lo que confirma que la
configuraci\'n entregada por \textsc{dhcp} permite alcanzar los servicios
internos:

\capturaAPA{fig:ping-gateway}{Prueba de conectividad desde el cliente hacia la puerta de enlace}{captura-ping-router-gateway.png}{Captura de pantalla de la terminal del cliente mostrando las respuestas de \texttt{ping} hacia la puerta de enlace \texttt{192.168.0.1} y hacia el servidor del proyecto.}

% ============================================================================
%  BLOQUE 4: SERVIDOR WEB CON PROXY INVERSO
% ============================================================================
\section{Servidor Web con proxy inverso}

El servidor web se implement\'o con Nginx, que cumple un doble rol: aloja
el sitio institucional del proyecto y act\'ua como proxy inverso
centralizado hacia los servicios internos que operan en puertos no
publicados, como el gestor de impresi\'on \textsc{cups}, el explorador de
archivos, el panel de gesti\'on de contenedores y el servidor de monitoreo
\cite{nginxDocs}.

\subsection{Arquitectura del servicio}

Todo el tr\'afico web entrante se protege mediante \textsc{https} con
certificados \textsc{tls} que incluyen nombres alternativos en formato
comodi\'n. La figura \ref{fig:arquitectura-web} ilustra el esquema
implementado.

\figuraAPA{fig:arquitectura-web}{Arquitectura del servidor web: Nginx como proxy inverso y servidor de contenido}{%
  \begin{tikzpicture}[
      node distance=1cm and 1.6cm,
      servicio/.style={draw, rounded corners, minimum width=2.4cm,
                       minimum height=0.8cm, align=center, font=\small},
      cliente/.style={draw, circle, minimum size=0.8cm, font=\small},
      linea/.style={-Stealth, thick}
    ]
    \node[servicio, fill=blue!10] (nginx) {Nginx\\{\footnotesize HTTP/HTTPS}};
    \node[servicio, right=2.2cm of nginx, yshift=2cm] (cups) {CUPS\\{\footnotesize :631}};
    \node[servicio, right=2.2cm of nginx, yshift=0.7cm] (files) {Filebrowser\\{\footnotesize :8082}};
    \node[servicio, right=2.2cm of nginx, yshift=-0.7cm] (portainer) {Portainer\\{\footnotesize :9000}};
    \node[servicio, right=2.2cm of nginx, yshift=-2cm] (zabbix) {Zabbix Web\\{\footnotesize :80}};
    \node[servicio, below=of nginx] (html) {/var/www/\\\footnotesize html};
    \node[cliente, left=of nginx] (browser) {Web};
    \draw[linea] (browser) -- node[above,font=\scriptsize]{443} (nginx);
    \draw[linea] (nginx) -- (cups);
    \draw[linea] (nginx) -- (files);
    \draw[linea] (nginx) -- (portainer);
    \draw[linea] (nginx) -- (zabbix);
    \draw[linea] (nginx) -- (html);
  \end{tikzpicture}%
}{Elaboraci\'on propia.}

La tabla \ref{tab:virtual-hosts} resume los \emph{virtual hosts}
configurados, sus dominios y su destino.

\begin{table}[H]
\centering
\caption{\emph{Virtual hosts} configurados en Nginx}
\label{tab:virtual-hosts}
\contenidoConFuente{%
  \begin{tablaAPA}
    \begin{tabular}{@{}p{4.6cm}p{4cm}p{5.2cm}@{}}
    \toprule
    \textbf{Subdominio} & \textbf{Destino} & \textbf{Funci\'on} \\
    \midrule
    \texttt{web.sudoers.lan} & \texttt{/var/www/html} & Sitio institucional del proyecto \\
    \texttt{print.sudoers.lan} & \texttt{127.0.0.1:631} & Proxy inverso al servicio de impresi\'on \\
    \texttt{files.sudoers.lan} & \texttt{127.0.0.1:8082} & Proxy inverso al gestor de archivos \\
    \texttt{portainer.sudoers.lan} & \texttt{127.0.0.1:9000} & Proxy inverso al panel de gesti\'on \\
    \texttt{zabbix.sudoers.lan} & \texttt{192.168.0.3:80} & Proxy inverso al monitoreo \\
    \texttt{\_} (por defecto) & Redirecci\'on a \textsc{https} & Acceso por IP o dominios no declarados \\
    \bottomrule
    \end{tabular}
  \end{tablaAPA}%
}{Elaboraci\'on propia.}
\end{table}

\subsection{Instalaci\'on y arranque del servicio}

\begin{lstlisting}[style=terminal]
$ sudo apt update
$ sudo apt install -y nginx openssl
$ systemctl enable --now nginx
\end{lstlisting}

El paquete \texttt{nginx} aporta el servidor web de arquitectura orientada a
eventos, adecuado para cargas concurrentes, y \texttt{openssl} permite
generar los certificados \textsc{x.509} requeridos para el protocolo
\textsc{https} \cite{opensslDocs}. Las directivas globales se ubican en
\texttt{/etc/nginx/nginx.conf} y cada \emph{virtual host} en
\texttt{/etc/nginx/sites-available/}.

\capturaAPA{fig:web-arranque}{Arranque del servicio Nginx en el servidor}{captura-web-arranque.png}{Captura de pantalla de la terminal mostrando la instalaci\'on de los paquetes y el arranque del servicio \texttt{nginx} con su estado activo.}

\subsection{Aplicaci\'on de \textsc{tls} y pruebas de acceso}

Tras emitir el certificado con \texttt{openssl} y asociarlo al
\emph{virtual host} se verific\'o la negociaci\'n segura y la publicaci\'on
del sitio a trav\'es del proxy inverso.

\capturaAPA{fig:web-cert}{Certificado \textsc{tls} emitido para el servidor web}{captura-web-cert.png}{Captura de pantalla de la salida de \texttt{openssl} mostrando los datos del certificado generado, incluido el nombre alternativo del dominio.}

\capturaAPA{fig:web-proxy}{Regla de proxy inverso configurada en Nginx}{captura-web-proxy.png}{Captura de pantalla de la configuraci\'n del \emph{virtual host} mostrando las directivas \texttt{proxy\_pass} y \texttt{proxy\_set\_header} que reencaminan la petici\'on al servicio interno.}

\capturaAPA{fig:web-sitio}{Sitio institucional publicado a trav\'es del proxy inverso}{bancosudoers.png}{Captura de pantalla del navegador mostrando el sitio web del proyecto cargado correctamente mediante el subdominio configurado.}

% ============================================================================
%  BLOQUE 5: CONTROLADOR DE DOMINIO SAMBA 4 ACTIVE DIRECTORY
% ============================================================================
\section{Controlador de dominio Samba 4 Active Directory}

El despliegue del dominio constituye el n\'ucleo de identidad de la
infraestructura. Samba 4 integra en un \'unico daemon los protocolos que
forman Active Directory: \textsc{ldap} para el directorio, Kerberos para la
autenticaci\'on y el \textsc{dns} propio del dominio
\cite{sambateam2024, rfc4511, rfc4120}.

\subsection{Preparaci\'on del sistema}

Un controlador de dominio es el punto de referencia de identidad de toda
la red, por lo que su nombre y su direcci\'on \textsc{ip} deben ser estables
y conocidos de antemano. Se fij\'o el nombre \texttt{dc1} y su nombre
completamente cualificado \texttt{dc1.sudoers.lan}:

\begin{lstlisting}[style=terminal]
$ hostnamectl set-hostname dc1
$ hostnamectl status | head -3
  Static hostname: dc1
  Operating System: Debian GNU/Linux 12 (bookworm)
\end{lstlisting}

\capturaAPA{fig:hostname}{Identidad del host \texttt{dc1} configurada en el servidor}{captura-hostname.png}{Captura de pantalla mostrando la salida de \texttt{hostnamectl status} con el nombre de host \texttt{dc1} y la distribuci\'on del sistema operativo.}

Un controlador de dominio no puede tener una direcci\'on variable, ya que
Kerberos y el \textsc{dns} lo referencian constantemente. Se garantiz\'o que
\texttt{dc1} resuelva a \texttt{192.168.0.2/24} en la interfaz cableada y se
registr\'o en \texttt{/etc/hosts}:

\begin{lstlisting}[style=terminal]
$ ip -4 a show eno1
    inet 192.168.0.2/24 brd 192.168.0.255 scope global eno1

$ grep -E 'dc1|sudoers' /etc/hosts
192.168.0.2   dc1.sudoers.lan   dc1
\end{lstlisting}

\capturaAPA{fig:ip-fija}{Direcci\'on \textsc{ip} fija \texttt{192.168.0.2/24} en la interfaz cableada}{captura-ip-fija.png}{Captura de pantalla de la salida de \texttt{ip -4 a show eno1} con la direcci\'on fija del servidor y la entrada de \texttt{/etc/hosts} que asocia \texttt{dc1.sudoers.lan} a esa direcci\'on.}

\subsection{Instalaci\'on de Samba 4 y promoci\'on del servidor}

Los paquetes del dominio se instalaron de forma nativa sobre el host, nunca
dentro de un contenedor: el controlador de Samba es un servicio sensible al
aislamiento de red y sistema de archivos que impone un contenedor
\cite{sambateam2024}.

\begin{lstlisting}[style=terminal]
$ apt update
$ apt install -y samba smbclient winbind krb5-user \
    libpam-winbind libnss-winbind dnsutils
\end{lstlisting}

Junto con Samba se habilitaron los servicios que sustentan el dominio:
el firewall \texttt{nftables}, el servidor de hora \texttt{chrony}
--necesario porque Kerberos rechazaTickets con un desfase de reloj superior
a cinco minutos-- y Docker para los servicios contenedorizados:

\begin{lstlisting}[style=terminal]
$ systemctl is-enabled samba-ad-dc nftables chrony docker
samba-ad-dc.service    enabled
nftables.service       enabled
chrony.service         enabled
docker.service         enabled
\end{lstlisting}

\capturaAPA{fig:servicios-host}{Servicios del host habilitados para el arranque}{captura-servicios-host.png}{Captura de pantalla de la salida de \texttt{systemctl is-enabled} con los servicios \texttt{samba-ad-dc}, \texttt{nftables}, \texttt{chrony} y \texttt{docker} en estado \texttt{enabled}.}

La promoci\'n a controlador de dominio se realiza aprovisionando el dominio
con \texttt{samba-tool}:

\begin{lstlisting}[style=terminal]
$ samba-tool domain provision \
    --use-rfc2307 --realm=SUDOERS.LAN \
    --domain=SUDOERS \
    --server-role=dc \
    --dns-backend=SAMBA_INTERNAL
\end{lstlisting}

La opci\'n \texttt{--use-rfc2307} activa la compatibilidad \textsc{posix}
de los grupos, de modo que la gesti\'on de permisos Unix se integra con la
del dominio.

\subsection{Configuraci\'on del \textsc{dns} del dominio}

El controlador de dominio administra la zona autoritativa interna. La
resoluci\'n se verific\'o con \texttt{dig} tanto para registros de servicio
como para las \textsc{srv} que publican Kerberos esta informaci\'on a
los clientes:

\begin{lstlisting}[style=terminal]
$ dig +short SRV _ldap._tcp.sudoers.lan
$ dig @dc1 sudoers.lan SOA
$ dig @dc1 dc1.sudoers.lan A
\end{lstlisting}

\capturaAPA{fig:zona-dns}{Zona autoritativa \texttt{sudoers.lan} administrada por el controlador de dominio}{captura-zona-dns.png}{Captura de pantalla del contenido de la zona DNS del dominio, con los registros de servicios y de los servidores que la integran.}

\capturaAPA{fig:dig-dc}{Resoluci\'n de nombres contra el controlador de dominio}{captura-dig-dc.png}{Captura de pantalla de la salida de \texttt{dig} con la resoluci\'n de los nombres del dominio contra el servidor \texttt{dc1}.}

\capturaAPA{fig:dig-srv}{Registros \textsc{srv} publicados por el dominio}{captura-dig-srv.png}{Captura de pantalla de la consulta de registros \textsc{srv} del dominio, que publishing las direcciones de los servicios Kerberos y \textsc{ldap}.}

\subsection{Configuraci\'on del servicio \textsc{dhcp} en el dominio}

El servicio \textsc{dhcp} se despliega en un contenedor con la opci\'n
\texttt{network\_mode: host}, condici\'n indispensable porque el protocolo
utiliza la direcci\'n de difusi\'n por \textsc{broadcast} en el puerto
\texttt{67/udp} y no funciona mapeado a trav\'es de un puente de red.

\capturaAPA{fig:kea-ss}{Estado del servicio \textsc{dhcp} del dominio en ejecuci\'on}{captura-kea-ss.png}{Captura de pantalla del estado del servicio de asignaci\'on de direcciones dentro del dominio, confirmar la entrega de direcciones a los clientes.}

\subsection{Administraci\'on de usuarios y grupos}

La gesti\'on de identidades se centraliza en el dominio, lo que evita
credenciales dispersas entre servicios. Las cuentas de administraci\'on se
asignan al grupo \texttt{admins} y las de prueba al grupo de oficina:

\begin{lstlisting}[style=terminal]
$ samba-tool user create joaquin --given-name="Joaquin" \
    --random-password
$ samba-tool group add admins
$ samba-tool group add users
$ samba-tool group add machines
\end{lstlisting}

\capturaAPA{fig:usuarios}{Usuarios creados en el dominio}{captura-usuarios.png}{Captura de pantalla de la salida de \texttt{samba-tool user list} mostrando las cuentas registradas en el dominio.}

\capturaAPA{fig:grupos}{Grupos de seguridad del dominio}{captura-grupos.png}{Captura de pantalla de la salida de \texttt{samba-tool group list} con los grupos definidos, incluidos los de administraci\'on y usuarios.}

\capturaAPA{fig:user-show}{Detalle de una cuenta de usuario del dominio}{captura-user-show.png}{Captura de pantalla de la informaci\'n detallada de una cuenta de usuario, incluyendo su pertenencia a los grupos del dominio.}

\subsection{Uni\'on de clientes y verificaci\'on del dominio}

La uni\'on de un cliente al dominio se verifica obteniendo un ticket de
Kerberos con \texttt{kinit} y list\'andolo con \texttt{klist}:

\begin{lstlisting}[style=terminal]
$ kinit joaquin
Password for joaquin@SUDOERS.LAN:
$ klist
Ticket cache: FILE:/tmp/krb5cc_0
Default principal: joaquin@SUDOERS.LAN
\end{lstlisting}

\capturaAPA{fig:kinit-cliente}{Obtenci\'on de ticket de Kerberos en el cliente}{captura-kinit-cliente.png}{Captura de pantalla de la terminal del cliente ejecutando \texttt{kinit} para obtener un ticket de Kerberos del dominio.}

\capturaAPA{fig:klist}{Listado de tickets de Kerberos del cliente}{captura-kinit-klist.png}{Captura de pantalla de la salida de \texttt{klist} mostrando la cach\'e de tickets y el principal por defecto del usuario autenticado.}

La resoluci\'n de grupos en el lado del cliente se verific\'o con
\texttt{getent}, que confirma la integraci\'on de \texttt{winbind} con la
resoluci\'on de nombres:

\capturaAPA{fig:getent}{Resoluci\'on de grupos del dominio en el cliente}{captura-getent.png}{Captura de pantalla de la salida de \texttt{getent group} mostrando los grupos del dominio resueltos correctamente en el cliente.}

El acceso a los recursos compartidos se verific\'o con \texttt{smbclient}
usando las credenciales del dominio:

\capturaAPA{fig:smbclient-ad}{Acceso a recursos compartidos con credenciales del dominio}{captura-smbclient.png}{Captura de pantalla de una sesi\'on \texttt{smbclient} autenticada con las credenciales del dominio, mostrando el acceso al recurso compartido.}

La informaci\'n general del dominio se consult\'o con las utilidades
propias de Samba:

\capturaAPA{fig:domain-info}{Informaci\'on del dominio y de la lista de equipos}{captura-domain-info.png}{Captura de pantalla de la salida de \texttt{samba-tool domain info} con los datos del dominio y su relaci\'on de confianza.}

\capturaAPA{fig:computer-list}{Equipos unidos al dominio}{captura-computer-list.png}{Captura de pantalla de la salida de \texttt{samba-tool computer list} mostrando los equipos que han unido al dominio.}

\subsection{Puertos publicados por el dominio}

El controlador de dominio expone un conjunto definido de puertos, cuya
disponibilidad se verific\'o desde el cliente:

\capturaAPA{fig:puertos-ad}{Verificaci\'on de los puertos de servicio del dominio}{captura-puertos-ad.png}{Captura de pantalla de la terminal del cliente mostrando los puertos del controlador de dominio que responden: DNS, Kerberos, \textsc{ldap} y SMB.}

% ============================================================================
%  BLOQUE 6: DNS DEL DOMINIO Y SERVIDOR WEB INTEGRADO (a8)
% ============================================================================
\section{DNS del dominio y servidor Web integrado}

La \'ultima etapa del proyecto consolid\'o en un \'unico dise\~no los dos
servicios que articulan toda la infraestructura: el \textsc{dns} autoritativo
del dominio \texttt{sudoers.lan}, servido por el propio controlador de
dominio, y el servidor web que publica los servicios de la intranet por
nombre. Este bloque describe la soluci\'n final, donde el \textsc{dns} y la
puerta web convergen sobre una \'unica direcci\'n de entrada \cite{rfc1034}.

\subsection{Arquitectura conjunta \textsc{dns} + \textsc{web}}

Como todos los servicios se ejecutan en el \'unico servidor
(\texttt{192.168.0.2}), el \textsc{dns} del dominio y la puerta web comparten
esa misma direcci\'n. La separaci\'on entre ambos es \textbf{l\'ogica} y no
f\'isica: el \textsc{dns} se atiende de forma nativa desde el daemon del
controlador de dominio, mientras que Nginx corre en un contenedor que
publica los puertos 80 y 443 \cite{tanenbaum2017}.

La tabla \ref{tab:mapa-nombres} resume los nombres gestionados por el
\textsc{dns} del dominio y el servicio al que dirige cada uno el proxy web.

\begin{table}[H]
\centering
\caption{Mapa de nombres del dominio y su backend en el proxy}
\label{tab:mapa-nombres}
\contenidoConFuente{%
  \begin{tablaAPA}
    \begin{tabular}{@{}p{5.2cm}p{4.4cm}p{4.2cm}@{}}
    \toprule
    \textbf{Nombre} & \textbf{Backend} & \textbf{Registro} \\
    \midrule
    \texttt{sudoers.lan}, \texttt{web}, \texttt{www} & Portal est\'atico (Nginx) & A \\
    \texttt{dc1.sudoers.lan} & Controlador de dominio (DNS, SMB) & A \\
    \texttt{print.sudoers.lan} & \texttt{tso-print:631} (CUPS) & A \\
    \texttt{portainer.sudoers.lan} & \texttt{tso-portainer:9000} & A \\
    \texttt{joaquin.sudoers.lan} & P\'agina del equipo (est\'atico) & A \\
    \texttt{david.sudoers.lan} & \texttt{192.168.0.51:8080} & A \\
    \texttt{nicolas.sudoers.lan} & \texttt{192.168.0.52:8080} & A \\
    \texttt{\_ldap.\_tcp.dc.\_msdcs} & \texttt{dc1}:389 (LDAP) & SRV \\
    \bottomrule
    \end{tabular}
  \end{tablaAPA}%
}{Elaboraci\'on propia.}
\end{table}

\subsection{El \textsc{dns} autoritativo servido por Samba}

Un dominio de Active Directory depende por completo del \textsc{dns}: los
clientes lo consultan para descubrir d\'onde reside cada servicio, tanto
mediante nombres comunes como mediante registros de descubrimiento
\textsc{srv} \cite{rfc2782, msft_activedirectory}. Por esta raz\'on, durante
el aprovisionamiento se indic\'o expl\'icitamente que el propio Samba
administrase la zona mediante el backend \textsc{samba\_internal}:

\begin{lstlisting}[style=terminal]
$ samba-tool domain provision \
    --realm=SUDOERS.LAN --domain=SUDOERS \
    --server-role=dc --dns-backend=SAMBA_INTERNAL \
    --use-rfc2307 --host-name=dc1 --host-ip=192.168.0.2 \
    --option="dns forwarder = 192.168.0.1"
\end{lstlisting}

Con este backend los registros se almacenan dentro del directorio y se
actualizan de forma din\'amica, sin necesidad de archivos de zona
manuales. La configuraci\'n resultante queda registrada en
\texttt{/etc/samba/smb.conf}:

\begin{lstlisting}[style=terminal]
$ grep -E 'server role|dns|realm' /etc/samba/smb.conf
[global]
    server role = active directory domain controller
    realm = SUDOERS.LAN
    dns forwarder = 192.168.0.1
    dns update command = /usr/sbin/samba_dnsupdate
\end{lstlisting}

\capturaAPA{fig:smb-dns}{Zona \texttt{sudoers.lan} y rol del controlador de dominio en \texttt{smb.conf}}{captura-smb-dns.png}{Captura de pantalla mostrando la salida del \texttt{grep} sobre \texttt{/etc/samba/smb.conf} con el rol \texttt{active directory domain controller}, el \textit{realm} \texttt{SUDOERS.LAN} y la directiva \texttt{dns forwarder = 192.168.0.1}.}

\subsection{Registros de la zona y de descubrimiento}

El contenido \'integral de la zona se consulta con \texttt{samba-tool dns},
que devuelve los registros \textsc{a} de todos los servicios publicados:

\begin{lstlisting}[style=terminal]
$ kinit administrator
$ samba-tool dns query dc1.sudoers.lan sudoers.lan @ ALL -k yes
Name=, Records=8, Children=0
  DC@SUDOERS.LAN: type=A, value=192.168.0.2
  web@SUDOERS.LAN: type=A, value=192.168.0.2
  print@SUDOERS.LAN: type=A, value=192.168.0.2
  portainer@SUDOERS.LAN: type=A, value=192.168.0.2
\end{lstlisting}

\capturaAPA{fig:zona-dns-a8}{Zona DNS completa del dominio consultada con \texttt{samba-tool dns}}{captura-zona-dns-a8.png}{Captura de pantalla mostrando la salida de \texttt{samba-tool dns query dc1.sudoers.lan sudoers.lan @ ALL} con los registros A de \texttt{dc1}, \texttt{web}, \texttt{print}, \texttt{portainer} y de los administradores, todos hacia \texttt{192.168.0.2}.}

La resoluci\'n de un nombre individual se verifica con \texttt{dig}, que
responde con la direcci\'n del proxy:

\begin{lstlisting}[style=terminal]
$ dig web.sudoers.lan
;; ANSWER SECTION:
web.sudoers.lan. 900 IN A 192.168.0.2
\end{lstlisting}

\capturaAPA{fig:dig-web}{Resoluci\'on de \texttt{web.sudoers.lan} por el \textsc{dns} del dominio}{captura-dig-web.png}{Captura de pantalla mostrando la salida de \texttt{dig web.sudoers.lan} con la respuesta \texttt{A 192.168.0.2}.}

Los registros \textsc{srv} de descubrimiento fueron generados
autom\'aticamente durante el aprovisionamiento y son los que permiten a los
clientes localizar el \textsc{ldap} del dominio:

\begin{lstlisting}[style=terminal]
$ dig _ldap._tcp.dc._msdcs.sudoers.lan SRV
;; ANSWER SECTION:
_ldap._tcp.dc._msdcs.sudoers.lan. 900 IN SRV 0 100 389 dc1.sudoers.lan.
\end{lstlisting}

\capturaAPA{fig:dig-srv-a8}{Registros \textsc{srv} de descubrimiento generados por el controlador de dominio}{captura-dig-srv-a8.png}{Captura de pantalla mostrando la salida de \texttt{dig \_ldap.\_tcp.dc.\_msdcs.sudoers.lan SRV} con el registro SRV apuntando al puerto 389 de \texttt{dc1.sudoers.lan}.}

\subsection{Registros de los equipos de administraci\'on}

Los nombres de los miembros del equipo tambi\'en resuelven hacia el proxy y
no hacia la direcci\'n de sus m\'aquinas. Gracias a esta indirecci\'on, un
backend puede cambiar de m\'aquina o de puerto sin que los nombres ni los
clientes se modifiquen. Los registros se crean con un script idempotente
que puede ejecutarse varias veces sin duplicarlos:

\begin{lstlisting}[style=terminal]
$ sudo kinit administrator
$ sudo bash dc/dns-records.sh
==> Servicios web alojados en maquinas de admins (via proxy) ...
  ok    david.sudoers.lan    ->  192.168.0.2  (ya existe)
  ok    nicolas.sudoers.lan  ->  192.168.0.2  (ya existe)
\end{lstlisting}

\capturaAPA{fig:dns-records}{Creaci\'on idempotente de los registros de los administradores}{captura-dns-records.png}{Captura de pantalla mostrando la salida de \texttt{sudo bash dc/dns-records.sh} con los registros \texttt{david} y \texttt{nicolas} apuntando a \texttt{192.168.0.2}.}

\subsection{Reenv\'io hacia Internet}

Para que los clientes resuelvan tambi\'en nombres externos, el controlador
act\'ua como \textit{forwarder} hacia el router de la red. De esta manera una
\'unica consulta al \textsc{dns} del dominio cubre tanto los nombres internos
como los de Internet, sin configurar un segundo servidor de nombres:

\begin{lstlisting}[style=terminal]
$ dig google.com @192.168.0.2 +short
142.250.78.78          # resolucion externa via el forwarder 192.168.0.1
\end{lstlisting}

\capturaAPA{fig:dig-externo}{Resoluci\'on de nombres externos a trav\'es del \textit{forwarder}}{captura-dig-externo.png}{Captura de pantalla mostrando la salida de \texttt{dig google.com @192.168.0.2} con una direcci\'on IP p\'ublica resuelta.}

\subsection{Integraci\'on con el servicio \textsc{dhcp}}

Para que una estaci\'n reci\'en conectada quede operativa sin intervenci\'on
manual, el servidor \textsc{dhcp} entrega junto con la direcci\'n las opciones
de red que apuntan al \textsc{dns} del dominio: la direcci\'n
\texttt{192.168.0.2}, el dominio \texttt{sudoers.lan}, la puerta de enlace
\texttt{192.168.0.1} y el servidor de hora. Toda m\'aquina de la red queda
as\'i capaz de resolver los nombres internos de forma transparente.

\subsection{Nginx como contenedor y proxy inverso}

El servidor web se despliega como servicio del \texttt{docker-compose} del
proyecto, a diferencia del controlador de dominio que corre de forma
nativa. La definici\'n del servicio establece los puertos publicados y los
vol\'umenes con el portal y los certificados:

\begin{lstlisting}[style=terminal]
web:
  build: ./services/web          # nginx:1.27-alpine
  container_name: tso-web
  ports:
    - "80:80"                    # redirige a HTTPS
    - "443:443"                  # HTTPS (termina TLS)
  volumes:
    - ./services/web/public:/usr/share/nginx/html:ro
    - ./services/web/tls:/etc/nginx/tls:ro      # CA + cert wildcard
  restart: always
\end{lstlisting}

\capturaAPA{fig:compose-web}{Servicio \texttt{web} definido en el \texttt{docker-compose.yml}}{captura-compose-web.png}{Captura de pantalla mostrando el bloque \texttt{web:} del \texttt{docker-compose.yml} con los puertos 80/443, los vol\'umenes del portal y de \textsc{tls}, y la directiva \texttt{restart: always}.}

Nginx enruta cada petici\'on por nombre y, adem\'as, redirige todo el
tr\'afico en claro hacia el canal cifrado:

\begin{lstlisting}[style=terminal]
server {
    listen 80;                                 # HTTP: solo redirige
    server_name sudoers.lan web.sudoers.lan www.sudoers.lan;
    return 301 https://$host$request_uri;
}

server {
    listen 443 ssl;                            # HTTPS: portal
    http2 on;
    server_name sudoers.lan web.sudoers.lan www.sudoers.lan;
    ssl_certificate     /etc/nginx/tls/server.crt;
    ssl_certificate_key /etc/nginx/tls/server.key;
    root  /usr/share/nginx/html;
    index index.html;
}
\end{lstlisting}

\capturaAPA{fig:nginx-conf}{Virtual host del portal y redirecci\'on de \textsc{http} a \textsc{https}}{captura-nginx-conf.png}{Captura de pantalla mostrando el contenido de la configuraci\'n de Nginx con el bloque \texttt{server} del portal y la directiva \texttt{return 301 https://\$host\$request\_uri}.}

\subsection{\textsc{https} con autoridad de certificaci\'on interna}

El portal se publica mediante \textsc{https} con un certificado
\textit{wildcard} que cubre todos los subdominios, firmado por una
autoridad de certificaci\'on interna generada con el script
\texttt{services/web/tls-gen.sh}:

\begin{lstlisting}[style=terminal]
$ sudo bash services/web/tls-gen.sh
=> Generando CA interna ...
=> Generando certificado para *.sudoers.lan ...
============================================
  CA:       ca.crt (instalar en los clientes para confiar)
  Server:   server.crt / server.key  (*.sudoers.lan, 825 dias)
============================================
$ docker compose up -d --build web
\end{lstlisting}

\capturaAPA{fig:tls-gen}{Generaci\'on de la autoridad interna y del certificado \textit{wildcard}}{captura-tls-gen.png}{Captura de pantalla mostrando la salida de \texttt{sudo bash services/web/tls-gen.sh} con la creaci\'on de la autoridad de certificaci\'on y del certificado \texttt{*.sudoers.lan}.}

El acceso al portal se verifica con \texttt{curl}, que devuelve el c\'odigo
de respuesta correcto:

\begin{lstlisting}[style=terminal]
$ curl -skI https://web.sudoers.lan
HTTP/1.1 200 OK
Server: nginx
\end{lstlisting}

\capturaAPA{fig:curl-web}{Acceso \textsc{https} al portal servido por Nginx}{captura-curl-web.png}{Captura de pantalla mostrando la salida de \texttt{curl -skI https://web.sudoers.lan} con el c\'odigo \texttt{200 OK} y el encabezado \texttt{Server: nginx}.}

\subsection{Proxy hacia los servicios de los administradores}

Los servicios alojados en los equipos de los administradores se exponen
tambi\'en a trav\'es del proxy. Sus nombres resuelven hacia
\texttt{192.168.0.2} y Nginx alcanza la m\'aquina correspondiente por la red
plana, de modo que el \'unico puerto visible desde la \textsc{lan} es el
443:

\begin{lstlisting}[style=terminal]
server {
    listen 443 ssl;
    server_name david.sudoers.lan;
    location / {
        resolver 127.0.0.11 valid=10s ipv6=off;
        set $david_upstream http://192.168.0.51:8080;
        proxy_pass $david_upstream;
        proxy_connect_timeout 5s;
    }
}
\end{lstlisting}

\begin{lstlisting}[style=terminal]
$ dig david.sudoers.lan +short          # el cliente resuelve el proxy
192.168.0.2
$ curl -sk https://david.sudoers.lan    # y el proxy responde la maquina
... (contenido del servicio de David) ...
\end{lstlisting}

\capturaAPA{fig:curl-david}{Servicio de la m\'aquina de David publicado a trav\'es del proxy}{captura-curl-david.png}{Captura de pantalla mostrando la resoluci\'on \texttt{dig david.sudoers.lan +short} con \texttt{192.168.0.2} y la respuesta de \texttt{curl -sk https://david.sudoers.lan} con el contenido del servicio del administrador.}

\subsection{Pruebas de integraci\'on desde un cliente}

La validaci\'on final se realiz\'o desde un equipo de la \textsc{lan}, que
recibi\'o su direcci\'n y su servidor \textsc{dns} por \textsc{dhcp}, y
confirm\'o que tanto la resoluci\'o de nombres como el acceso web funcionan
sin configuraci\'on manual:

\begin{lstlisting}[style=terminal]
$ dig web.sudoers.lan +short             # desde el cliente
192.168.0.2
$ curl -sk https://web.sudoers.lan | head -5
<!DOCTYPE html>
<html lang="es">
... (portal de la intranet) ...
\end{lstlisting}

\capturaAPA{fig:cliente-dig}{Resoluci\'on de nombres desde un cliente de la red}{captura-cliente-dig.png}{Captura de pantalla mostrando la salida de \texttt{dig web.sudoers.lan} ejecutada en un cliente, respondiendo \texttt{192.168.0.2}.}

\capturaAPA{fig:cliente-web}{Acceso \textsc{https} al portal desde el navegador del cliente}{captura-cliente-web.png}{Captura de pantalla del navegador de un cliente mostrando el portal de la intranet cargado por \texttt{https://web.sudoers.lan} con el candado de confianza de la autoridad interna.}

La misma validaci\'on se repiti\'o para el portal, la consola de impresi\'on,
el panel de gesti\'on de contenedores y las p\'aginas de los miembros, lo que
confirm\'o que el \textsc{dns} y el proxy web operan como un \'unico sistema
de nombres y de publicaci\'on.

% ============================================================================
%  BLOQUE 7: MEDIDAS DE SECURIZACI\'ON
% ============================================================================
\section{Medidas de seguridad aplicadas}

La seguridad de la infraestructura se abord\'o de forma transversal,
aplicando controles espec\'ificos en cada capa de la arquitectura descrita
en el cap\'itulo anterior \cite{stewart2017}.

\subsection{Seguridad perimetral y control de acceso}

El firewall \texttt{nftables} opera con pol\'itica \textsc{drop} por
defecto, de modo que todo el tr\'afico no permitido se descarta
expl\'icitamente. El acceso al puerto \texttt{22/tcp} se restringe a los
rango de direcciones reservadas para los administradores de TI, de modo
que los puestos operativos no pueden alcanzar la consola del servidor:

\begin{lstlisting}[style=terminal]
# Politica por defecto: descartar todo
$ nft add table inet filter
$ nft add chain inet filter input '{ type filter hook input priority 0; policy drop; }'

# Permitir administracion solo desde el rango de TI
$ nft add rule inet filter input ip saddr 192.168.0.50-192.168.0.52 \
    tcp dport 22 ct state new accept
\end{lstlisting}

\subsection{Aislamiento de servicios internos}

Los servicios que no deben exponerse a la red se mantienen en puertos
internos y se publican \'unicamente a trav\'es del proxy inverso Nginx. De
esta forma, la base de datos \textsc{postgresql} y el panel de gesti\'on de
contenedores no son alcanzables de forma directa desde la \textsc{lan}.

\subsection{Autenticaci\'on y control de acceso por rol}

La gesti\'on centralizada de identidades en el dominio permite asignar
permisos a los recursos compartidos por grupo de usuarios, en lugar de
otorgar acceso individual a cada estaci\'on. La autorizaci\'on del correo
se restricti\'a a los usuarios autenticados mediante \textsc{sasl}, y el
env\'io an\'onimo queda deshabilitado con la opci\'on
\texttt{noanonymous}.

\subsection{Cifrado del tr\'afico}

Las comunicaciones sensibles se protegen con \textsc{tls} mediante
certificados emitidos para los dominios internos, lo que impide el
tr\'ansito de credenciales en texto claro tanto en el env\'io de correo
como en el acceso web.

% ============================================================================
%  BLOQUE 7: PRUEBAS Y VALIDACI\'ON
% ============================================================================
\section{Pruebas funcionales y validaci\'on de la infraestructura}

La validaci\'on final se realiz\'o comprobando la resolutividad de los
nombres del dominio, la obtenci\'on de una direcci\'on v\'alida por
\textsc{dhcp} y el acceso a los servicios publicados a trav\'es del proxy
inverso. Las pruebas ejecutadas desde el cliente fueron:

\begin{lstlisting}[style=terminal]
$ ping -c 4 192.168.0.2
$ dig @192.168.0.2 web.sudoers.lan
$ curl -I https://web.sudoers.lan
$ kinit joaquin
\end{lstlisting}

La obtenci\'n de la direcci\'n por \textsc{dhcp} se valid\'o mediante la
salida de \texttt{dhclient}, en la que el cliente renueva la concesion
previa y recibe una direcci\'n dentro del \'ambito din\'amico, junto con la
m\'ascara, la puerta de enlace y el servidor \textsc{dns} del proyecto.
Esta renovaci\'n result\'o especialmente significativa porque la direcci\'n
anterior proven\'ia del servicio \textsc{dhcp} del router, y su
sustituci\'n confirm\'o que el servidor del dominio asume \'unicamente la
totalidad de las concessions de la red.

Los resultados obtenidos confirman que los tres servicios operan de forma
conjunta: el cliente obtiene su direcci\'n, resuelve los nombres del dominio
contra el controlador de dominio y accede a los servicios publicados a
trav\'es del proxy inverso, validando la integraci\'on integral de la
infraestructura.


% --- 5. DESAFÍOS DE IMPLEMENTACIÓN Y SOLUCIONES APLICADAS -------------------
\caratulaCapitulo{5}{DESAF\'IOS DE IMPLEMENTACI\'ON Y SOLUCIONES APLICADAS}
% ============================================================================
%  5. DESAF\'IOS DE IMPLEMENTACI\'ON Y SOLUCIONES APLICADAS
%  ----------------------------------------------------------------------------
%  Los desafíos descritos corresponden a dificultades reales
%  documentadas durante la ejecución de las actividades a4, a5 y a7.
% ============================================================================

\section{Desaf\'ios de implementaci\'on y soluciones aplicadas}

La puesta en marcha de la infraestructura plante\'o dificultades reales
que no se anticipate durante la planificaci\'on inicial. Este cap\'itulo
recoge los cinco desaf\'ios m\'as significativos, describe su sintomatizaci\'on
y expone la soluci\'n adoptada por el grupo, incluyendo las verificaciones
que permitieron confirmar la resoluci\'on.

\subsection{Conflicto entre dos servidores \textsc{dhcp} activos}

El primer obstacle se manifest\'o al conectar el primer cliente a la red:
obten\'ia una direcci\'n de \texttt{192.168.0.1}, correspondiente al router,
en lugar de recibirla del servidor Kea del proyecto. La causa era que el
router ten\'ia su propio servicio \textsc{dhcp} activo, de modo que tanto
\'este como Kea respond\'ian a los paquetes \textsc{discover} con ofertas
distintas. El cliente aceptaba la primera oferta recibida, lo que
provocaba una configuraci\'n de red inconsistente y, en consecuencia,
direcciones duplicadas entre equipos \cite{rfc2131}.

Para resolverlo se accedi\'o a la interfaz de administraci\'on del router y
se deshabilit\'o el servicio \textsc{dhcp} en la \textsc{lan}, manteniendo
\'unicamente la configuraci\'on de la puerta de enlace y las reglas de
red. La figura \ref{fig:desafio-dhcp} documenta el estado final de dicho
par\'ametro.

\capturaAPA{fig:desafio-dhcp}{Desactivaci\'on del servicio \textsc{dhcp} del router para eliminar el conflicto con Kea}{captura-router-disable.png}{Captura de pantalla de la configuraci\'on \textsc{lan} del router TP-Link donde la opci\'on \textsc{dhcp server} figura en estado deshabilitado, resolviendo la Respuesta duplicada de los dos servidores \textsc{dhcp}.}

\subsection{Verificaci\'on previa de la conectividad del cliente}

Tras la desactivaci\'on del \textsc{dhcp} del router se comprob\'o que
ning\'un servicio de la red hubiera quedado inaccesible como efecto
secundario. Esta verificaci\'on result\'o relevante porque el router
concentra la puerta de enlace y la salida a Internet: si alg\'un equipo
hab\'ia quedado sin ruta por defecto, el cambio habr\'ia afectado a todos
los equipos que dependieran de ella. La figura \ref{fig:desafio-ping}
muestra la verificaci\'on realizada.

\capturaAPA{fig:desafio-ping}{Verificaci\'on de la conectividad del cliente hacia la puerta de enlace tras el cambio}{captura-ping-router-gateway.png}{Captura de pantalla de la terminal del cliente con las respuestas de \texttt{ping} hacia la puerta de enlace \texttt{192.168.0.1}, confirmando que el equipo mantiene su conectividad tras la intervenci\'on.}

\subsection{Resoluci\'on err\'onea del \textit{hostname} a \textit{loopback}}

Al instalar el controlador de dominio \texttt{dc1} se observ\'o que las
operaciones de uni\'on al dominio fallaban con errores de acceso al
servidor Kerberos. El origen estaba en la resoluci\'on del nombre del host:
si \texttt{dc1} se resuelve a \texttt{127.0.0.1}, el servidor se comunica
consigo mismo a trav\'es del \textit{loopback} y los clientes de la red
nunca alcanzan al controlador de dominio \cite{sambateam2024}.

La soluci\'n consisti\'o en garantizar que el nombre \texttt{dc1} resuelva a
la direcci\'n \texttt{192.168.0.2/24} de la interfaz cableada \texttt{eno1}
y registrar esa correspondencia de forma explícita en \texttt{/etc/hosts},
de modo que la resoluci\'o local tenga prioridad sobre cualquier otra
entrada.

\capturaAPA{fig:desafio-loopback}{Correspondencia entre el nombre del host y su direcci\'n de red en \texttt{/etc/hosts}}{captura-ip-fija.png}{Captura de pantalla de la salida de \texttt{ip -4 a show eno1} junto con la entrada de \texttt{/etc/hosts} que asocia \texttt{dc1.sudoers.lan} a la direcci\'n de la red local y no al \textit{loopback}.}

\subsection{Desfase de reloj y autenticaci\'on Kerberos}

Las lecturas de tickets de Kerberos fallaban de forma intermitente en los
clientes, con el error de versi\'on de ticket incorrecta. La causa no era la
configuraci\'on del dominio sino el reloj del sistema: Kerberos exige que
el desfase entre el cliente y el servidor sea inferior a cinco minutos
para aceptar el ticket \cite{rfc4120}.

La soluci\'on fue instalar y habilitar el servicio \texttt{chrony} en el
host del controlador de dominio, de modo que la hora del servidor se
sincronizara de forma continua con los servidores de referencia. Tras
activar la sincronizaci\'on, la obtenci\'on de tickets se complet\'o de
forma satisfactoria.

\capturaAPA{fig:desafio-kerberos}{Servicios de sincronizaci\'on horaria habilitados en el host del dominio}{captura-servicios-host.png}{Captura de pantalla de la salida de \texttt{systemctl is-enabled} mostrando \texttt{chrony} junto con \texttt{samba-ad-dc}, \texttt{nftables} y \texttt{docker} habilitados en el arranque del servidor.}

\subsection{Concesi\'on \textsc{dhcp} heredada del router}

Tras resolver el conflicto de servidores \textsc{dhcp}, los equipos que ya
ten\'ian una direcci\'n asignada por el router la conservaban, ya que la
concesi\'on anterior segu\'ia siendo v\'alida. Estos clientes no
actualizaban su configuraci\'on de red y, en consecuencia, persist\'an sin
poder resolver los nombres del dominio \cite{droms1997}.

La soluci\'on consisti\'o en forzar la renovaci\'on de la concesion desde
el cliente, de modo que \'este la solicite al servidor Kea y reciba una
direcci\'n dentro del \'ambito din\'amico junto con la puerta de enlace y el
servidor \textsc{dns} correctos. Esta renovaci\'n tambi\'n sirvi\'o como
prueba de que el servidor \textsc{dhcp} del proyecto asum\'ia \'unicamente
la totalidad de las concesiones de la red.

\capturaAPA{fig:desafio-dhclient}{Renovaci\'on forzada de la concesion \textsc{dhcp} en el cliente}{captura-dhclient.png}{Captura de pantalla de la terminal del cliente mostrando la renovaci\'on de la direcci\'n y la asignaci\'n de una direcci\'n dentro del \'ambito din\'amico definido por Kea.}

\subsection{Servicios que requieren red \textit{host} en contenedores}

Por \'ultimo, el despliegue del servicio \textsc{dhcp} dentro de Docker
present\'o un inconveniente de conectividad: el protocolo utiliza la
direcci\'on de difusi\'on por \textsc{broadcast} en el puerto
\texttt{67/udp}, y los contenedores con red \textit{bridge} no reenv\'ian
ese tipo de tr\'afico. Como resultado, los clientes no recib\'ian oferta
alguna.

La soluci\'on fue configurar el contenedor con la opci\'n
\texttt{network\_mode: host}, que descarta el aislamiento de red y permite
al servicio escuchar directamente en la interfaz del host. Esta decisi\'on
se tom\'o de forma consciente, asumiendo que el servicio queda expuesto a
la \textsc{lan} y que por ello debe protegerse mediante las reglas de
firewall del host \cite{stewart2017}.

\subsection{Advertencia de seguridad del navegador por autoridad de certificaci\'on no reconocida}

Al publicar el portal por \textsc{https} con el certificado
\textit{wildcard} \texttt{*.sudoers.lan}, los navegadores de los clientes
mostraban una advertencia de que la conexi\'on no era privada. El servidor
respond\'a correctamente y el cifrado es real, pero el certificado lo firma
una autoridad de certificaci\'on interna generada por el propio grupo y, por
tanto, no figura en la lista de autoridades en las que el navegador conf\'ia
por defecto. Para el usuario esto se manifiesta como un riesgo inexistente
que induce a ignorar la advertencia.

La soluci\'on consisti\'o en distribuir el certificado de la autoridad
(\texttt{ca.crt}) e instalarlo como entidad de confianza en el sistema de
cada estaci\'on. Una vez instalada, el candado aparece en la barra de
direcciones sin advertencias y el cifrado queda validado de extremo a
extremo. La figura \ref{fig:desafio-tls} documenta el acceso al portal una
vez resuelta la cadena de confianza.

\capturaAPA{fig:desafio-tls}{Acceso al portal tras instalar la autoridad de certificaci\'on interna en el cliente}{captura-cliente-web.png}{Captura de pantalla del navegador de un cliente mostrando el portal de la intranet con el candado de confianza de la autoridad interna, sin la advertencia de certificado no reconocido.}


% --- 6. CONCLUSIONES --------------------------------------------------------
\caratulaCapitulo{6}{CONCLUSIONES}
% ============================================================================
%  6. CONCLUSIONES
%  ----------------------------------------------------------------------------
%  Cierra el informe cubriendo las cuatro actividades del proyecto.
% ============================================================================

\section{Conclusiones}

\subsection{Resultados obtenidos}

El objetivo central del trabajo se cumpli\'o: se implement\'o y valid\'o una
arquitectura de red empresarial funcional y segura sobre GNU/Linux, en la
que los servicios operan de forma integrada y no como piezas aisladas.

En el \'ambito del correo electr\'onico, el servicio qued\'o operativo con
Postfix como agente de transferencia y Dovecot como servidor de acceso a
buzones, con autenticaci\'on \textsc{sasl} y protecci\'on \textsc{starttls}.
El env\'io autenticado, la recepci\'n en formato \texttt{maildir} y el acceso
desde cliente de escritorio se verificaron satisfactory. En el \'ambito de
archivos, los tres protocolos operan sobre un mismo directorio compartido,
de modo que un archivo creado a trav\'es de \textsc{smb}, \textsc{nfs} o
\textsc{http} es inmediatamente accesible desde los otros, lo que evita la
duplicaci\'on de la informaci\'on almacenada.

La asignaci\'on din\'amica de direcciones qued\'o centralizada en el servidor
\textsc{kea}, tras resolverse el conflicto con el servicio \textsc{dhcp} del
router. Los clientes reciben su direcci\'on, la puerta de enlace y el
servidor \textsc{dns} sin intervenci\'on manual, y las direcciones de
infraestructura permanecen estables mediante reservas. El servidor web con
proxy inverso public\'o los servicios internos por nombre, con
redirecci\'on autom\'atica de \textsc{http} a \textsc{https} y una autoridad
de certificaci\'on interna que cubre todos los subdominios.

El despliegue del controlador de dominio Samba 4 fue el hito que permiti\'o
consolidar el resto. Centraliz\'o la identidad, la autenticaci\'on
\textsc{kerberos} y el \textsc{ldap}, y aport\'o el \textsc{dns} autoritativo
del dominio con sus registros de descubrimiento \textsc{srv}. Gracias a
ello, un equipo cliente se une al dominio, resuelve los nombres internos y
accede a los servicios publicados, todo ello a trav\'es de una \'unica
puerta de entrada cifrada, sin configurar direcciones ni puertos en el
cliente.

\subsection{Importancia del proyecto en la formaci\'on profesional}

El proyecto result\'o especialmente formativo porque no se limit\'o a la
configuraci\'on aislada de herramientas, sino que exigi\'o comprender c\'omo
se relacionan los servicios entre s\'i. Un error en la configuraci\'on del
servidor de nombres se manifiesta como un fallo aparentemente ajeno, como
la imposibilidad de autenticar un usuario en el correo; un par\'ametro
mal definido en \textsc{dhcp} impide a los clientes alcanzar al servidor que
precisamente les fue entregado. Comprender esas relaciones es lo que
distingue a un administrador de red de quien simplemente ejecuta
instrucciones.

El entorno de trabajo result\'o especialmente valioso por ser
realista. Las restricciones encontradas --un \textsc{dhcp} activo en el
router que compet\'ia con el del proyecto, la necesidad de sincronizar el
reloj para que \textsc{kerberos} aceptara los tickets, la confianza de la
autoridad de certificaci\'on en cada estaci\'on-- no se deducen de la
teor\'ia: se descubren diagnosticando. Ese constituye el principal
aprendizaje del ejercicio profesional.

\subsection{Aprendizajes adquiridos}

Se adquirieron competencias en cuatro \'ambitos. En el dise\~no de redes,
se aprendi\'o a segmentar una red plana en zonas con pol\'iticas de acceso
distintas, y a justificar cada decisi\'on t\'ecnica con argumentos en lugar
de aplicarla por costumbre.

En la administraci\'on de servicios, se aprendi\'o a configurar correo
electr\'onico, compartici\'on de archivos, asignaci\'on din\'amica de
direcciones, publicaci\'on web y un directorio corporativo,
comprendiendo el rol espec\'ifico que desempeña cada uno dentro del conjunto y
las dependencias que los unen. Se comprendi\'o tambi\'en la diferencia
pr\'actica entre ejecutar un servicio de forma nativa y desplegarlo en un
contenedor, y por qu\'e el controlador de dominio deb\'ia ejecutarse de la
primera forma \cite{sambateam2024}.

En el \'ambito de la seguridad, se aprendi\'o a aplicar la pol\'itica de
rechazo por defecto, a restringir el acceso administrativo a rangos de
direcciones determinados, a aislar los servicios sensibles tras un \'unico
proxy y a confiar el tr\'afico mediante \textsc{tls} con una autoridad
interna \cite{stewart2017}.

En el diagnostico, el principal aprendizaje fue el m\'etodo: leer el
registro del servicio, contrastar la configuraci\'on efectiva con la
esperada y verificar con la herramienta de l\'inea de comandos antes de
modificar nada. Esta rutina result\'o m\'as valiosa que el conocimiento
memorizado de una opci\'n concreta, porque se transfiere a cualquier
servicio que no se haya visto antes \cite{salz2020}.

\subsection{Recomendaciones para futuras implementaciones}

A partir de la experiencia desarrollada se formulan las siguientes
recomendaciones.

\begin{itemize}
  \item \textbf{Automatizar el aprovisionamiento.} Aunque la
        instalaci\'on se realiz\'o de forma manual, el car\'acter
        declarativo de \textsc{samba-tool} y de los archivos de configuraci\'on
        en \textsc{json} de \textsc{kea} permite escribir guiones que
        reconstruyan el dominio completo de forma idempotente, lo que
        resulta indispensable ante la necesidad de recuperar un servidor
        tras un fallo.
  \item \textbf{Automatizar la renovaci\'on de certificados.} El
        certificado \textit{wildcard} generado tiene una vigencia
        acotada. Un proceso programado que lo reemita y lo recargue
        evitar\'ia las interrupciones de servicio por expiraci\'on.
  \item \textbf{Centralizar los registros.} La resoluci\'on de nombres
        externa depende hoy de un \textit{forwarder} hacia el router, un
        punto \'unico de falla. Instalar un \textsc{dns} recursivo local o
        configurar un \textsc{dns} de respaldo eliminar\'ia esa
        dependencia.
  \item \textbf{Implementar el monitoreo de forma efectiva.} El servidor de
        monitoreo est\'a desplegado, pero a\'un no se han configurado
        alertas que notifiquen la ca\'ida de un servicio. Conectar las
        alertas con el servidor de correo cerrar\'ia el ciclo de
        notificaci\'on.
  \item \textbf{Definir y probar la copia de seguridad.} El proceso de
        respaldo debe especificar qu\'e se copia, con qu\'e frecuencia y c\'omo
        se verifica la restauraci\'on. Una copia nunca comprobada no
        constituye un respaldo.
  \item \textbf{Documentar cada cambio.} Mantener un registro de las
        modificaciones aplicadas a la configuraci\'on permiti\'o resolver
        los problemas del cap\'itulo anterior. Convertir esa pr\'actica en
        un procedimiento formalizado es el siguiente paso natural.
\end{itemize}


% --- Referencias (APA) ------------------------------------------------------
\newpage
\encabezadoCapitulo{REFERENCIAS}{}
\addcontentsline{toc}{part}{REFERENCIAS}
\bibliographystyle{apacite}
\nocite{*}
\bibliography{formato/referencias}

\end{document}
